\documentclass[12pt,a4paper]{article}

\usepackage[utf8]{inputenc}
\usepackage{geometry}
\usepackage{amsmath}
\usepackage{booktabs}
\usepackage{array}
\usepackage{longtable}
\usepackage{enumitem}
\usepackage{hyperref}
\usepackage{float}

\geometry{margin=1in}

\title{
    \textbf{CSAI 302 -- Advanced Database Systems}\\
    \large Lab 2: Buffer Pool Management, Storage Models, and Page Encryption
}
\author{}
\date{Spring 2026}

\begin{document}

\maketitle

\section{Lab Objectives}

The main objectives of this lab are:

\begin{enumerate}
    \item Implement a Buffer Pool Manager.
    \item Implement and compare the following replacement policies:
    \begin{itemize}
        \item LRU
        \item Clock
        \item LRU-K
        \item 2Q
    \end{itemize}
    \item Benchmark the replacement policies using different buffer pool sizes.
    \item Calculate hit and miss ratios.
    \item Implement row-oriented storage (NSM).
    \item Implement column-oriented storage (DSM).
    \item Convert a row-oriented representation into a column-oriented representation.
    \item Compare row-oriented and column-oriented storage.
    \item Encrypt and decrypt database pages.
    \item Benchmark encrypted and non-encrypted storage.
\end{enumerate}

\section{Part 1: Buffer Pool Manager}

\subsection{1.1 What is a Buffer Pool?}

A database stores a large amount of data on disk. However, disk access is much slower than accessing data in main memory (RAM).

Therefore, a database does not read every page from disk every time it is requested.

Instead, frequently used pages are temporarily kept in memory inside a structure called the \textbf{Buffer Pool}.

A simple analogy is a student studying from books.

Suppose the books are stored in a library:

\begin{itemize}
    \item The library represents the disk.
    \item The student's desk represents RAM.
    \item A book represents a database page.
    \item A location on the desk represents a buffer frame.
\end{itemize}

If the student already has the required book on the desk, there is no need to go to the library.

Similarly, if a database page is already in the buffer pool, the database can use it immediately.

\subsection{1.2 Page and Frame}

A \textbf{page} is a fixed-size unit of data managed by the database.

For example, a database may use pages of:

\[
4KB,\quad 8KB,\quad 16KB
\]

depending on the system.

A \textbf{frame} is a slot in the buffer pool that can contain one page.

For example, if the buffer pool has three frames:

\begin{verbatim}
Frame 0     Frame 1     Frame 2
+-------+   +-------+   +-------+
|Page 2 |   |Page 5 |   |Page 8 |
+-------+   +-------+   +-------+
\end{verbatim}

The database has three pages currently resident in memory.

\subsection{1.3 Page Table}

The Buffer Pool Manager needs to know where every resident page is located.

It therefore maintains a mapping:

\[
PageID \rightarrow FrameID
\]

For example:

\begin{verbatim}
Page 2 -> Frame 0
Page 5 -> Frame 1
Page 8 -> Frame 2
\end{verbatim}

This structure is called the \textbf{Page Table}.

\subsection{1.4 Page Hit}

Suppose the query requests Page 5.

The Buffer Pool Manager checks the page table:

\begin{verbatim}
Page 5 -> Frame 1
\end{verbatim}

Therefore, Page 5 is already in memory.

This is called a:

\[
\boxed{\text{Hit}}
\]

No disk read is required.

\subsection{1.5 Page Miss}

Suppose the query requests Page 10.

The page table does not contain Page 10.

Therefore:

\[
\boxed{\text{Miss}}
\]

The Buffer Pool Manager must obtain Page 10 from disk.

If there is an empty frame, Page 10 is loaded directly into that frame.

If all frames are occupied, the Buffer Pool Manager must choose a victim page using a replacement policy.

\subsection{1.6 Dirty Pages}

A page may be modified while it is in memory.

For example, suppose disk contains:

\begin{verbatim}
Page 5:
GPA = 3.5
\end{verbatim}

After an update in memory:

\begin{verbatim}
Page 5:
GPA = 3.9
\end{verbatim}

The in-memory version is now different from the disk version.

Therefore, the page is marked:

\[
dirty = True
\]

If a dirty page is selected as a victim, it must be written back to disk before its frame can be reused.

A clean page does not need to be written back because its disk version is already up to date.

\subsection{1.7 Pin Count}

A page may currently be used by a query.

Such a page should not be evicted while it is being used.

The Buffer Pool Manager can therefore maintain:

\[
pin\_count
\]

For example:

\begin{verbatim}
Page 5:
pin_count = 2
\end{verbatim}

This means that two users or operations are currently using the page.

Every \emph{fetch} increments the pin count and every \emph{unpin} decrements it. A page with:

\[
pin\_count > 0
\]

is not eligible for replacement. Only when the pin count returns to $0$ may the replacement policy choose it.

\subsection{1.8 Complete Page Fetch Flow}

When a query requests Page X, the Buffer Pool Manager performs the following steps:

\begin{enumerate}
    \item Check the page table.
    \item If Page X exists, return it immediately (Hit).
    \item Otherwise, record a Miss.
    \item Check for a free frame.
    \item If a free frame exists, use it.
    \item Otherwise, ask the replacement policy to select a victim.
    \item Make sure the victim is not pinned.
    \item If the victim is dirty, write it to disk.
    \item Read Page X from disk.
    \item Place Page X into the selected frame.
    \item Update the page table.
    \item Increment the pin count and return Page X to the caller.
\end{enumerate}

The complete process is:

\begin{verbatim}
            Request Page X
                  |
      Is Page X in the page table?
          /                \
        Yes                 No
         |                   |
        HIT                MISS
         |                   |
         |            Free frame exists?
         |              /          \
         |            Yes           No
         |             |             |
         |             |    Replacement policy
         |             |    selects a victim
         |             |             |
         |             |       Victim dirty?
         |             |        /        \
         |             |      Yes         No
         |             |       |           |
         |             |  Write victim     |
         |             |  to disk          |
         |             |       |           |
         |             +-------+-----------+
         |                     |
         |              Read Page X from disk
         |              into the frame and
         |              update the page table
         |                     |
         +---------------------+
                  |
          pin_count += 1
                  |
             Return Page X
\end{verbatim}

\subsection{1.9 Hit Ratio and Miss Ratio}

The effectiveness of a buffer pool is measured by how often a requested page is already in memory:

\[
Hit\ Ratio =
\frac{Hits}{Hits+Misses}
\qquad
Miss\ Ratio =
\frac{Misses}{Hits+Misses}
\]

Since every access is either a hit or a miss:

\[
Hit\ Ratio + Miss\ Ratio = 1
\]

A higher hit ratio means fewer disk reads. Two factors affect it:

\begin{itemize}
    \item \textbf{Buffer pool size}: more frames can keep more pages in memory, so the hit ratio usually increases with the buffer size.
    \item \textbf{Replacement policy}: when the buffer is full, a better choice of victim keeps the pages that will be needed again.
\end{itemize}

However, a larger buffer helps only up to a limit. The first access to every distinct page must be read from disk; this is called a \textbf{compulsory miss}. For example, a workload of 18 requests to 8 distinct pages always has at least 8 misses, so once the buffer can hold all 8 pages, adding more frames does not improve the hit ratio.

\section{Part 2: Replacement Policies}

When the buffer pool is full, the database must decide which page should be removed.

This is the responsibility of the replacement policy. Only unpinned pages can be chosen.

In this lab, four policies are implemented:

\begin{enumerate}
    \item LRU
    \item Clock
    \item LRU-K
    \item 2Q
\end{enumerate}

\subsection{2.1 LRU}

LRU stands for:

\[
\boxed{\text{Least Recently Used}}
\]

The idea is simple:

\begin{quote}
Remove the page that has not been used for the longest time.
\end{quote}

For example, assume three frames and the access sequence:

\[
A \rightarrow B \rightarrow C \rightarrow A
\]

After accessing A again, the order from least to most recently used is:

\begin{verbatim}
B (oldest)  ->  C  ->  A (most recent)
\end{verbatim}

If a new page D arrives, B is selected as the victim.

\subsubsection{LRU Example}

Consider a buffer pool with three frames and the reference string:

\[
2,3,2,1,5,3,4,3
\]

The execution is:

\begin{longtable}{c|c|c|c}
\toprule
Step & Page & Result & Buffer Pool \\
\midrule
1 & 2 & Miss & [2,-,-] \\
2 & 3 & Miss & [2,3,-] \\
3 & 2 & Hit & [2,3,-] \\
4 & 1 & Miss & [2,3,1] \\
5 & 5 & Miss & [2,5,1] \\
6 & 3 & Miss & [3,5,1] \\
7 & 4 & Miss & [3,5,4] \\
8 & 3 & Hit & [3,5,4] \\
\bottomrule
\end{longtable}

There are:

\[
Hits = 2
\]

and:

\[
Misses = 6
\]

Therefore:

\[
Hit\ Ratio =
\frac{Hits}{Hits+Misses}
=
\frac{2}{8}
=
0.25
\]

Thus:

\[
\boxed{Hit\ Ratio=25\%}
\]

and:

\[
\boxed{Miss\ Ratio=75\%}
\]

\subsubsection{Weakness: Sequential Flooding}

Suppose pages A and B are used constantly (a ``hot'' set), and then a query scans pages
$1, 2, 3$ exactly once. With three frames:

\begin{verbatim}
A, B, A, B   -> [A, B, -]
scan 1       -> [A, B, 1]
scan 2       -> [B, 1, 2]    (A evicted)
scan 3       -> [1, 2, 3]    (B evicted)
A, B         -> Miss, Miss
\end{verbatim}

Pages that will never be used again pushed out the useful pages. This is called
\textbf{sequential flooding}, and it motivates LRU-K and 2Q.

\subsection{2.2 Clock Replacement}

Clock is an approximation of LRU.

Instead of maintaining an exact ordering of all pages, Clock uses:

\begin{itemize}
    \item A circular list of frames.
    \item A clock hand.
    \item A reference bit for each frame.
\end{itemize}

Each frame has:

\[
reference\_bit \in \{0,1\}
\]

When a page is accessed, its reference bit is set to 1.

The clock hand scans frames.

\begin{itemize}
    \item If reference bit = 1, set it to 0 and give the page a second chance.
    \item If reference bit = 0, the page can be selected as a victim.
\end{itemize}

\subsubsection{Clock Example}

Suppose:

\begin{verbatim}
Frame 0 -> A, ref=1
Frame 1 -> B, ref=0
Frame 2 -> C, ref=1
\end{verbatim}

The clock hand points to A.

We need a victim.

First:

\begin{verbatim}
A: ref=1
\end{verbatim}

Therefore, A receives a second chance:

\begin{verbatim}
A: ref 1 -> 0
\end{verbatim}

The hand moves to B.

B has:

\begin{verbatim}
B: ref=0
\end{verbatim}

Therefore, B is selected as the victim.

The important idea is:

\begin{quote}
A recently used page receives a second chance before being evicted.
\end{quote}

Clock is cheaper than LRU because an access only sets a bit, instead of reordering a list.
Because it approximates LRU, it also suffers from sequential flooding.

\subsection{2.3 LRU-K}

LRU-K extends the idea of LRU by considering multiple historical accesses.

The value $K$ specifies how many accesses are considered.

For example:

\begin{itemize}
    \item LRU-2 considers the last two accesses.
    \item LRU-3 considers the last three accesses.
\end{itemize}

The value $K$ is not the number of pages.

It is the amount of access history used for each page.

\subsubsection{Access History Example}

Consider:

\begin{verbatim}
Step:  1 2 3 4 5 6
Page:  A B A C B A
\end{verbatim}

The access history is:

\begin{verbatim}
A -> 1,3,6
B -> 2,5
C -> 4
\end{verbatim}

For LRU-2:

\begin{verbatim}
A -> last two accesses = 3,6
B -> last two accesses = 2,5
C -> only one access
\end{verbatim}

For a page with at least $K$ accesses, the backward $K$-distance at current time $t$ can be expressed as:

\[
D_K(p)=t-H_K(p)
\]

where $H_K(p)$ is the $K$-th most recent access time of page $p$.

A larger backward $K$-distance means that the relevant historical access is older.

Pages with fewer than $K$ recorded accesses have an undefined $K$-th historical access; canonical LRU-K treats their backward $K$-distance as $+\infty$, so they are evicted first. When several pages have $+\infty$ distance, a tie-breaking rule is needed; in this lab, the page with the oldest access is evicted.

This is why LRU-K resists sequential flooding: pages read only once by a scan have infinite distance and are evicted before pages that were used repeatedly.

\subsubsection{Simple LRU-2 Example}

Suppose the history is:

\begin{verbatim}
Step:  1 2 3 4 5 6 7
Page:  A B C A B C A
\end{verbatim}

At step 7:

\begin{verbatim}
A -> 1,4,7
B -> 2,5
C -> 3,6
\end{verbatim}

For LRU-2, use the second-most-recent access:

\begin{verbatim}
A -> 4
B -> 2
C -> 3
\end{verbatim}

At time $t=7$:

\[
D_2(A)=7-4=3
\]

\[
D_2(B)=7-2=5
\]

\[
D_2(C)=7-3=4
\]

The largest distance is for B, so B is selected as the victim.

\subsection{2.4 2Q}

2Q stands for:

\[
\boxed{\text{Two-Queue}}
\]

The goal is to prevent one-time sequential scans from destroying useful pages in the buffer pool.

2Q distinguishes between:

\begin{itemize}
    \item Pages that have been seen recently but only once.
    \item Pages that have demonstrated repeated use.
\end{itemize}

The two main queues are:

\begin{verbatim}
A1 -> recently seen/new pages       (FIFO)
Am -> pages with repeated use       (LRU)
\end{verbatim}

A new page enters A1. If it is accessed again while still in A1, it is promoted to Am.
Victims are taken from A1 first, so pages seen only once cannot push out pages in Am.

A commonly used full 2Q design also maintains an $A1\text{-out}$ ghost queue containing identifiers (not data) of pages recently evicted from $A1$; a page found there on its next access goes directly to Am.

\subsubsection{2Q Example}

Assume four frames and the sequence:

\[
A,B,A,C,D,E
\]

Initially:

\begin{verbatim}
A1 = []
Am = []
\end{verbatim}

After accessing A:

\begin{verbatim}
A1 = [A]
Am = []
\end{verbatim}

After B:

\begin{verbatim}
A1 = [A,B]
Am = []
\end{verbatim}

A is accessed again.

A has now demonstrated repeated use, so it is promoted:

\begin{verbatim}
A1 = [B]
Am = [A]
\end{verbatim}

After C and D, all four frames are full:

\begin{verbatim}
A1 = [B,C,D]
Am = [A]
\end{verbatim}

When E arrives, the victim is taken from A1 (the oldest, B), not from Am:

\begin{verbatim}
A1 = [C,D,E]
Am = [A]
\end{verbatim}

The key idea is:

\begin{quote}
A page that was accessed only once should not immediately become more valuable than a page that has demonstrated repeated use.
\end{quote}

\subsection{2.5 Summary of Policies}

\begin{center}
\begin{tabular}{p{0.12\textwidth}|p{0.45\textwidth}|p{0.3\textwidth}}
\toprule
Policy & Victim & Sequential flooding \\
\midrule
LRU & Least recently used page & Suffers \\
Clock & First page with reference bit 0 & Suffers (approximates LRU) \\
LRU-K & Largest backward $K$-distance & Resists \\
2Q & Oldest page in A1, then LRU in Am & Resists \\
\bottomrule
\end{tabular}
\end{center}

\subsection{2.6 Effect of the Access Pattern}

Which policy performs best depends on the access pattern:

\begin{itemize}
    \item \textbf{Uniform random}: no policy can predict the next request, so all policies perform similarly.
    \item \textbf{Skewed (80--20)}: a small set of pages receives most requests; policies that remember repeated use (LRU-K, 2Q) keep those pages.
    \item \textbf{Hot set + sequential scan}: LRU and Clock suffer from sequential flooding, while LRU-K and 2Q protect the hot set.
\end{itemize}

\section{Part 3: Row-Oriented Storage (NSM)}

NSM stands for:

\[
\boxed{\text{N-ary Storage Model}}
\]

In row-oriented storage, attributes belonging to the same tuple are stored together.

Consider:

\begin{center}
\begin{tabular}{c|c|c|c}
\toprule
ID & Name & Age & GPA \\
\midrule
1 & Ali & 20 & 3.5 \\
2 & Sara & 21 & 3.8 \\
3 & Omar & 20 & 3.2 \\
\bottomrule
\end{tabular}
\end{center}

A row-oriented representation is:

\begin{verbatim}
Row 1 = [1, Ali, 20, 3.5]
Row 2 = [2, Sara, 21, 3.8]
Row 3 = [3, Omar, 20, 3.2]
\end{verbatim}

Inside a page, complete rows are placed one after another:

\begin{verbatim}
Page 0: [header][1,Ali,20,3.5][2,Sara,21,3.8][3,Omar,20,3.2]...
\end{verbatim}

\subsection{Advantages of Row Storage}

Row storage is useful when a query needs most or all attributes of a tuple.

For example:

\begin{verbatim}
SELECT *
FROM Students
WHERE ID = 2;
\end{verbatim}

The system needs the complete tuple:

\begin{verbatim}
[2, Sara, 21, 3.8]
\end{verbatim}

All of it is in a single page, so one page read is enough.

Row storage is therefore commonly associated with transactional workloads (OLTP).

\section{Part 4: Column-Oriented Storage (DSM)}

DSM stands for:

\[
\boxed{\text{Decomposition Storage Model}}
\]

In column-oriented storage, values belonging to the same attribute are stored together.

The same table becomes:

\begin{verbatim}
ID   = [1, 2, 3]

Name = [Ali, Sara, Omar]

Age  = [20, 21, 20]

GPA  = [3.5, 3.8, 3.2]
\end{verbatim}

Each column is stored in its own set of pages:

\begin{verbatim}
ID pages:   [header][1][2][3]...
Name pages: [header][Ali][Sara][Omar]...
GPA pages:  [header][3.5][3.8][3.2]...
\end{verbatim}

To return a complete row, the system must read one page from every column and combine the values. This is called \textbf{tuple reconstruction}.

\section{Part 5: Row Store to Column Store Conversion}

The conversion process reads each row and appends each attribute to the corresponding column.

Suppose:

\begin{verbatim}
Rows:

[1, Ali, 20, 3.5]
[2, Sara, 21, 3.8]
[3, Omar, 20, 3.2]
\end{verbatim}

After conversion:

\begin{verbatim}
ID   = [1, 2, 3]
Name = [Ali, Sara, Omar]
Age  = [20, 21, 20]
GPA  = [3.5, 3.8, 3.2]
\end{verbatim}

The same data is stored in both layouts; only the physical organization changes. After conversion, the results of any query must be identical in both stores.

\section{Part 6: Comparing Row and Column Storage}

The better layout depends on the query pattern.

\subsection{Full Row Query}

Example:

\begin{verbatim}
SELECT *
FROM Students
WHERE ID = 2;
\end{verbatim}

This type of access needs many attributes from one tuple and is naturally suitable for row-oriented storage.

\subsection{Single Column Aggregation}

Example:

\begin{verbatim}
SELECT AVG(GPA)
FROM Students;
\end{verbatim}

Only the GPA column is required.

A column-oriented representation can scan:

\begin{verbatim}
GPA = [3.5, 3.8, 3.2, ...]
\end{verbatim}

without needing to process unrelated columns.

\subsection{Group By Query}

Example:

\begin{verbatim}
SELECT Age, AVG(GPA)
FROM Students
GROUP BY Age;
\end{verbatim}

The important columns are:

\begin{verbatim}
Age
GPA
\end{verbatim}

A column-oriented layout is therefore well suited to this analytical access pattern.

\subsection{Why Columns Read Fewer Pages}

Assume 4KB pages ($4096$ bytes) and rows of $(ID, Name, GPA)$ with sizes $4 + 20 + 4 = 28$ bytes:

\[
\text{Rows per page (NSM)} \approx \frac{4096}{28} \approx 146
\qquad
\text{GPA values per page (DSM)} \approx \frac{4096}{4} = 1024
\]

For $200{,}000$ students, computing \texttt{AVG(GPA)} reads approximately:

\[
\frac{200{,}000}{146} \approx 1370 \text{ pages (NSM)}
\qquad
\frac{200{,}000}{1024} \approx 196 \text{ pages (DSM)}
\]

The column store reads about $7\times$ fewer pages. For a full-row lookup, the opposite happens: the row store reads 1 page, while the column store reads 1 page per column.

\subsection{Storage Comparison}

\begin{center}
\begin{tabular}{p{0.25\textwidth}|p{0.3\textwidth}|p{0.3\textwidth}}
\toprule
 & Row-Oriented (NSM) & Column-Oriented (DSM) \\
\midrule
Basic unit & Row/Tuple & Column/Attribute \\
Full-row access & Good & Requires reconstruction \\
Few-column analytical queries & Less efficient & Efficient \\
OLTP & Commonly suitable & Usually less natural \\
OLAP & Less suitable for wide rows & Commonly suitable \\
Aggregation & May read unnecessary attributes & Reads selected columns \\
\bottomrule
\end{tabular}
\end{center}

\section{Part 7: Encrypting and Decrypting Data Pages}

Database pages may contain sensitive information.

For example:

\begin{verbatim}
Page 7

Student ID = 10
Name = Ali
GPA = 3.8
\end{verbatim}

Storing the page as plaintext means that someone who obtains the database files may potentially read the contents directly.

Encryption transforms plaintext into ciphertext.

\[
Plaintext
\xrightarrow{Encryption}
Ciphertext
\]

Decryption performs the reverse operation:

\[
Ciphertext
\xrightarrow{Decryption}
Plaintext
\]

The important architecture is:

\begin{verbatim}
             BUFFER POOL
                  |
             Plain Page
                  |
               Encrypt
                  |
                  v
                DISK
            Encrypted Page
\end{verbatim}

When the page is read:

\begin{verbatim}
                DISK
            Encrypted Page
                  |
               Decrypt
                  |
                  v
             BUFFER POOL
               Plain Page
\end{verbatim}

\subsection{7.1 AES-GCM}

AES-GCM is a mode of authenticated encryption.

It provides:

\begin{enumerate}
    \item Confidentiality.
    \item Integrity/authentication.
\end{enumerate}

Confidentiality means that the encrypted contents cannot be directly understood without the key.

Integrity/authentication means that unauthorized modification can be detected.

Conceptually:

\[
Ciphertext,\ Tag
=
AES\text{-}GCM\_Encrypt(Key, Nonce, Plaintext)
\]

Decryption verifies the authentication tag:

\[
Plaintext
=
AES\text{-}GCM\_Decrypt(Key, Nonce, Ciphertext, Tag)
\]

If the ciphertext has been modified, authentication should fail.

In this lab, a 256-bit AES key is used.

\subsection{7.2 Nonce}

AES-GCM requires a nonce.

A nonce must never be reused with the same encryption key; reusing it breaks the security of AES-GCM.

A common practical choice is a fresh random 12-byte nonce for every encryption operation.

Therefore, two encryptions can have:

\begin{verbatim}
Page 7 -> Nonce N1
Page 8 -> Nonce N2
Page 7 again -> Nonce N3
\end{verbatim}

A fresh nonce also means that writing the same page contents twice produces two different ciphertexts.

The nonce does not need to be secret, but it must be stored with the encrypted data because it is needed for decryption.

\subsection{7.3 Authentication Tag}

AES-GCM produces an authentication tag.

Conceptually, the encrypted page can contain:

\begin{verbatim}
Nonce (12 bytes)
+
Ciphertext (same size as the page)
+
Authentication Tag (16 bytes)
\end{verbatim}

Therefore, each encrypted page on disk is 28 bytes larger than the plain page.

During decryption, the tag is checked.

If someone modifies the ciphertext, the verification fails.

Therefore:

\begin{verbatim}
Correct Ciphertext
      |
      v
Authentication OK
      |
      v
Decrypt
\end{verbatim}

while:

\begin{verbatim}
Modified Ciphertext
      |
      v
Authentication FAILED
      |
      v
Reject
\end{verbatim}

\subsection{7.4 Page ID as Associated Data}

The page identifier can optionally be supplied as Additional Authenticated Data (AAD).

For example:

\[
AAD = PageID
\]

The AAD is not encrypted, but it is included in the tag computation.

This allows the authentication check to bind the ciphertext to its expected page identifier.

If ciphertext belonging to Page 8 is presented as Page 7, authentication fails.

This protects against page-swapping or page-mix-up scenarios.

\section{Part 8: Integrating Encryption with the Buffer Pool}

Encryption should normally happen at the disk boundary.

Suppose Page 5 is dirty:

\begin{verbatim}
Buffer Pool:
Page 5
GPA = 3.9
dirty = True
\end{verbatim}

When Page 5 is evicted:

\begin{verbatim}
Page 5
   |
Encrypt using AES-GCM
   |
   v
Encrypted Page 5
   |
   v
Disk
\end{verbatim}

When Page 5 is requested later:

\begin{verbatim}
Disk
 |
Encrypted Page 5
 |
Decrypt
 |
Plain Page 5
 |
Buffer Pool
\end{verbatim}

Therefore, the buffer pool can operate on plaintext pages while persistent storage contains encrypted pages.

An important consequence is that encryption and decryption happen \textbf{only on a miss or when a dirty page is written back}. A hit does not touch the disk, so it pays no encryption cost. The higher the hit ratio, the smaller the encryption overhead.

\section{Part 9: Important Concepts to Remember}

\subsection*{Buffer Pool}

The buffer pool is a memory area used to cache database pages and reduce disk I/O.

\subsection*{Page}

A page is a fixed-size unit of database storage.

\subsection*{Frame}

A frame is a slot in the buffer pool that holds one page.

\subsection*{Hit}

A requested page is already in the buffer pool.

\subsection*{Miss}

A requested page is not in the buffer pool and must be fetched from disk.

\subsection*{Pin Count}

The number of operations currently using a page; a pinned page cannot be evicted.

\subsection*{Dirty Page}

A page modified in memory that must be written back to disk before eviction.

\subsection*{LRU}

Evicts the page whose most recent access is the oldest.

\subsection*{Clock}

Uses reference bits and a circular clock hand to provide pages with a second chance.

\subsection*{LRU-K}

Uses the last $K$ accesses of each page to make a replacement decision.

\subsection*{2Q}

Separates recently seen pages from pages that have demonstrated repeated use.

\subsection*{Sequential Flooding}

A one-time scan evicts useful pages from an LRU-based buffer pool.

\subsection*{NSM}

Stores complete tuples/rows together.

\subsection*{DSM}

Stores values of the same attribute/column together.

\subsection*{Tuple Reconstruction}

Combining values from several columns to rebuild a complete row in a column store.

\subsection*{Encryption}

Transforms plaintext into ciphertext to protect data at rest.

\subsection*{Decryption}

Transforms ciphertext back into plaintext.

\subsection*{AES-GCM}

An authenticated encryption mode that provides confidentiality and integrity/authentication.

\section{Part 10: Final System Architecture}

The complete lab can be viewed as the following architecture:

\begin{verbatim}
                      DATABASE QUERY
                            |
             Row Store (NSM) / Column Store (DSM)
                            |
                            v
                   Buffer Pool Manager
                            |
                 +----------+----------+
                 |                     |
                HIT                   MISS
                 |                     |
                 |           Replacement Policy
                 |        (LRU | Clock | LRU-K | 2Q)
                 |                     |
                 |                  Victim
                 |                     |
                 |          Dirty? -> AES-GCM Encrypt
                 |                     |       |
                 |                     |       v
                 |                     |     DISK
                 |                     |       |
                 |              AES-GCM Decrypt (requested page)
                 |                     |
                 +----------+----------+
                            |
                        Plain Page
                            |
                          Query
\end{verbatim}

The complete system therefore demonstrates several important database system concepts:

\begin{enumerate}
    \item Memory management through a buffer pool.
    \item Cache replacement through multiple replacement policies.
    \item Performance evaluation using hit/miss ratios.
    \item Row-oriented and column-oriented physical data organization.
    \item Conversion between physical storage layouts.
    \item Protection of persistent database pages through encryption.
    \item Performance trade-offs introduced by encryption.
\end{enumerate}

\section{Conclusion}

This lab demonstrates how a database system manages the movement of pages between disk and memory.

The Buffer Pool Manager avoids unnecessary disk accesses by caching frequently used pages in RAM. When the buffer pool becomes full, a replacement policy such as LRU, Clock, LRU-K, or 2Q decides which page can be removed.

The lab also demonstrates two physical storage organizations. NSM stores complete tuples together and is useful for full-row access, while DSM stores values of the same attribute together and is useful for analytical queries that access a subset of columns.

Finally, page encryption protects data stored on disk. AES-GCM provides both confidentiality and authentication, at the cost of extra work on every disk read and write; because hits do not touch the disk, a high hit ratio keeps this cost small.

This lab therefore illustrates an important database-system principle:

\[
\boxed{
\text{Performance, Storage Organization, Memory Management, and Security are interconnected.}
}
\]

\end{document}
